% ==================================================================
% id:       comb-011
% topic:    קומבינטוריקה
% subtopic: המקדם הבינומי
% solution: full
% verified: no
% ==================================================================
היעזרו בבינום של ניוטון, $(a+b)^n=\sum_{k=0}^n\binom nka^kb^{n-k}$, והראו שלכל קבוצה סופית $\Omega$, מספר תתי-הקבוצות של $\Omega$ בגודל זוגי שווה למספר תתי-הקבוצות בגודל אי-זוגי. (רמז: סמנו $|\Omega|=n$ ובחרו את $a,b$ להיות\dots)
\begin{solbox}
\unverified
נסמן $|\Omega|=n$, ונציב בבינום של ניוטון $a=-1$, $b=1$:
\[
0=\big((-1)+1\big)^n=\sum_{k=0}^n\binom nk(-1)^k\,1^{n-k}=\binom n0-\binom n1+\binom n2-\binom n3+\dots+(-1)^n\binom nn.
\]
נעביר את המחוברים עם הסימן השלילי לאגף השני ונקבל
\[
\ans{\binom n1+\binom n3+\dots=\binom n0+\binom n2+\dots}
\]
(המחובר האחרון בכל אגף הוא $\binom nn$ או $\binom n{n-1}$, לפי הזוגיות של $n$). האגף הימני סופר את תתי-הקבוצות של $\Omega$ בגודל זוגי, והשמאלי את אלה בגודל אי-זוגי.
\end{solbox}
